% Section 2.9, from lectures/L02/appendix/b_exercises.md.

\section{Exercises}
\label{c2:sec:exercises}

\begin{aside}
\textbf{How to check your work.} Every exercise below that asks you to write a VHDL module comes
with a self-checking testbench under \file{lectures/L02/exercises/}. Write your module in its
directory \file{exercises/<module>/}, with the entity name and the \textbf{port order} the exercise
states, and then run it with GHDL \textemdash{} see \secref{c2:sec:testbenches} for the three
commands.

No FPGA board is needed for any exercise. The Quartus synthesis and board programming steps are
demonstrated during the session; your job afterwards is to get the VHDL right, and the testbench is
how you confirm it.
\end{aside}

% ------------------------------------------------------------------------------------------
\exerciseset{Karnaugh maps}

\exercise{Three variables}{Circuit}
\label{c2:ex:kmap3}
Derive a minimised equation for \code{X} from the Karnaugh map below, then realise the
corresponding gate network and simulate it in CircuitVerse:

\begin{narrowtable}{cc}
  \thead{ABC} & \thead{X} \\
  \midrule
  000 & 1 \\
  001 & 1 \\
  010 & 0 \\
  011 & 0 \\
  100 & 1 \\
  101 & 1 \\
  110 & 0 \\
  111 & 0 \\
\end{narrowtable}

\note[Hint]Put \code{AB} down the rows in Gray code order (\code{00, 01, 11, 10}) and \code{C}
across the columns, exactly as in the worked example in \secref{c2:sec:kmap}, before you start
grouping \code{1}s.

\exercise{Four variables}{Circuit}
\label{c2:ex:kmap4}
Derive a minimised equation for \code{X} from the Karnaugh map below, then realise the
corresponding gate network and simulate it in CircuitVerse:

\begin{narrowtable}{cc}
  \thead{ABCD} & \thead{X} \\
  \midrule
  0000 & 0 \\
  0001 & 1 \\
  0010 & 0 \\
  0011 & 1 \\
  0100 & 0 \\
  0101 & 0 \\
  0110 & 0 \\
  0111 & 0 \\
  1000 & 0 \\
  1001 & 1 \\
  1010 & 0 \\
  1011 & 1 \\
  1100 & 1 \\
  1101 & 0 \\
  1110 & 1 \\
  1111 & 0 \\
\end{narrowtable}

\note[Hint]With four variables you put two of them (in Gray code) on each axis, so that you get a
4x4 grid. Groups can still wrap around every edge of the grid, not just left\--right; remember to
check the top and bottom edges too.

% ------------------------------------------------------------------------------------------
\exerciseset{Circuits in VHDL}

\exercise{\code{xyz_logic}: three outputs from four inputs}{VHDL}
\label{c2:ex:xyz}
A gate network has four inputs, \code{ABCD}, and three outputs, \code{XYZ}. It is given by the
truth table below:

\begin{narrowtable}{cc}
  \thead{ABCD} & \thead{XYZ} \\
  \midrule
  0000 & 001 \\
  0001 & 010 \\
  0010 & 001 \\
  0011 & 011 \\
  0100 & 100 \\
  0101 & 110 \\
  0110 & 100 \\
  0111 & 111 \\
  1000 & 101 \\
  1001 & 110 \\
  1010 & 101 \\
  1011 & 111 \\
  1100 & 000 \\
  1101 & 010 \\
  1110 & 000 \\
  1111 & 011 \\
\end{narrowtable}

\xpart{a} Treat \code{X}, \code{Y} and \code{Z} as three separate single-output functions of
\code{ABCD}, and derive a minimised equation for each by way of a Karnaugh map.

\xpart{b} Realise the resulting gate network and simulate it in CircuitVerse; check all three
outputs against the truth table above.

\xpart{c} Implement the design in VHDL, in a module of its own with four \code{std_logic} inputs
and three \code{std_logic} outputs. \Secref{c1:sec:vhdl} and \secref{c2:sec:tovhdl} show the
pattern for turning a derived equation into concurrent VHDL assignments.

\xpart{d} Run the testbench and confirm that all 16 rows pass.

If a row fails, the assertion message names the input combination:
\begin{itemize}
  \item Go back to the Karnaugh map for whichever of \code{X}, \code{Y}, \code{Z} is wrong.
  \item A single wrong output on a single row is nearly always a mis-grouped \code{1} in a map.
\end{itemize}

\note[Hint]Two of the three outputs here turn out to depend on only one or two of the four inputs
once minimised. Do not assume every output needs all four.

\modulefig[0.52]{lectures/L02/appendix/images/xyz_logic.png}

\begin{selfcheck}
\textbf{Self-check.} Name your VHDL entity \code{xyz_logic}, with the ports \code{a}, \code{b},
\code{c}, \code{d} (in) and \code{x}, \code{y}, \code{z} (out), declared in that order; its
testbench is in \file{exercises/xyz_logic/} and checks all 16 rows of the truth table above.
\end{selfcheck}

% ------------------------------------------------------------------------------------------
\exerciseset{Multiplexers}

\exercise{The 4-to-1 mux by hand}{Circuit}
\label{c2:ex:mux4hand}
A 4-to-1 multiplexer has four data inputs, \code{A}--\code{D}, two select bits, \code{S[1:0]}, and
one output, \code{X}.

Let the selector name its data input directly: \code{S[1:0] = "00"} selects \code{A}, \code{"01"}
selects \code{B}, \code{"10"} selects \code{C} and \code{"11"} selects \code{D}. That is the mapping
Exercise~\ref{c2:ex:mux4} implements.

\xpart{a} Write out the multiplexer's truth table (four rows, one per select combination).

\xpart{b} Derive the sum-of-products equation for \code{X} in terms of \code{A}--\code{D} and
\code{S[1:0]}. \Secref{c2:sec:mux} does the same for the 8-to-1 case; halve that here.

\xpart{c} Build and simulate the network in CircuitVerse, and confirm that \code{X} follows each
data input in turn as you change \code{S[1:0]}.

Keep your answer: Exercise~\ref{c2:ex:mux4} below implements exactly this circuit in VHDL, once
\secref{c2:sec:process} has given you the \code{process} and \code{case} statements it needs.

\exercise{\code{mux2} with \code{if}/\code{else}}{VHDL}
\label{c2:ex:mux2}
\Secref{c2:sec:process} writes the 2:1 multiplexer with a \code{case} statement. Write the same
circuit once more with an \code{if}/\code{else} instead, and convince yourself that the two
describe identical hardware.

Write an entity named \code{mux2} with:

\begin{booktable}{@{}p{3.6em}p{3.4em}p{6.4em}L@{}}
  \thead{Port} & \thead{Dir.} & \thead{Type} & \thead{Description} \\
  \midrule
  \code{d0} & in & \code{std_logic} & Data input. \\
  \code{d1} & in & \code{std_logic} & Data input. \\
  \code{sel} & in & \code{std_logic} & Selector. \\
  \code{x} & out & \code{std_logic} & The selected data input. \\
\end{booktable}

The selector connects through the data input whose number it names:
\begin{itemize}
  \item \code{sel = '0'} selects \code{d0}.
  \item \code{sel = '1'} selects \code{d1}.
  \item Note that this is the reverse of the \code{case} example in \secref{c2:sec:process}, which
    selects its \emph{second} input \code{b} at \code{'0'}. Which input a select value names is a
    convention, not a rule, so read it off the specification every time rather than assuming.
\end{itemize}

Implement the multiplexer with a process-based architecture:
\begin{itemize}
  \item Include \code{d0}, \code{d1} and \code{sel} in the sensitivity list.
  \item Use an \code{if}/\code{else} statement on \code{sel}.
  \item Assign \code{x} in every branch.
  \item Do not use:
  \begin{itemize}
    \item A conditional concurrent assignment (\code{x <= d1 when sel = '1' else d0;}).
    \item A \code{with...select} statement: a concurrent form of \code{case} that the book does not
      cover, mentioned only so that you do not reach for it if you have met it elsewhere.
  \end{itemize}
\end{itemize}

\modulefig[0.5]{lectures/L02/appendix/images/mux2.png}

\begin{selfcheck}
\textbf{Self-check.} Name your entity \code{mux2}, with the ports \code{d0}, \code{d1}, \code{sel}
(in) and \code{x} (out), declared in that order; its testbench is in \file{exercises/mux2/} and
checks all eight combinations of the three inputs.
\end{selfcheck}

\exercise{\code{mux4}}{VHDL}
\label{c2:ex:mux4}
Extend Exercise~\ref{c2:ex:mux2} to a 4:1 multiplexer named \code{mux4}.

This is the circuit you derived, drew and simulated by hand in Exercise~\ref{c2:ex:mux4hand}.
Derive the \code{case} branches afresh yourself for four inputs and two select bits rather than
copying the 8-to-1 example from \secref{c2:sec:process}.

The entity must have:

\begin{booktable}{@{}p{3.6em}p{3.4em}p{10.8em}L@{}}
  \thead{Port} & \thead{Dir.} & \thead{Type} & \thead{Description} \\
  \midrule
  \code{d0} & in & \code{std_logic} & Data input. \\
  \code{d1} & in & \code{std_logic} & Data input. \\
  \code{d2} & in & \code{std_logic} & Data input. \\
  \code{d3} & in & \code{std_logic} & Data input. \\
  \code{sel} & in & \code{std_logic_vector(1 downto 0)} & Selector. \\
  \code{x} & out & \code{std_logic} & The selected data input. \\
\end{booktable}

As in Exercise~\ref{c2:ex:mux2}, the selector names its data input directly: \code{"00"} selects
\code{d0}, \code{"01"} selects \code{d1}, \code{"10"} selects \code{d2} and \code{"11"} selects
\code{d3}. Again this is the reverse mapping from \secref{c2:sec:process}'s worked 8-to-1 example,
where \code{"000"} connects \code{inputs(7)} through. Derive your \code{case} branches from the four
rows above, not from that listing.

Implement the multiplexer with a process:
\begin{itemize}
  \item Include every data input and \code{sel} in the sensitivity list.
  \item Use a \code{case} statement on \code{sel}, one branch per select value above.
  \item Include a \code{when others} branch, and choose a safe output value for select inputs
    containing values such as \code{'X'}, \code{'U'} or \code{'Z'}.
  \item Do not replace the \code{case} statement with an \code{if}/\code{elsif} chain.
\end{itemize}

\modulefig[0.56]{lectures/L02/appendix/images/mux4.png}

\begin{selfcheck}
\textbf{Self-check.} Name your entity \code{mux4}, with the ports \code{d0}, \code{d1}, \code{d2},
\code{d3}, \code{sel} (\code{std_logic_vector(1 downto 0)}) and \code{x} (out), declared in that
order; its testbench is in \file{exercises/mux4/} and sweeps all four select values against all
sixteen data input combinations.
\end{selfcheck}

% ------------------------------------------------------------------------------------------
\exerciseset{Submodules and 7-segment displays}

Exercise~\ref{c2:ex:display} and \ref{c2:ex:hexdisplay} build one design in two halves: a module
that drives a single 7-segment display, and a top level that uses two copies of it to show one
hexadecimal digit on each. Do them in order: Exercise~\ref{c2:ex:hexdisplay} instantiates what you
write in Exercise~\ref{c2:ex:display}.

A 7-segment display is seven LEDs, named \code{a} to \code{g}, arranged around a figure eight.
Light the right subset and you get a digit. This is the layout, and it is what you need in order to
work out the codes below:

\begin{textdrawing}
      aaaa
     f    b
     f    b
      gggg
     e    c
     e    c
      dddd
\end{textdrawing}

So \code{1} lights \code{b} and \code{c} (the two on the right), and \code{7} lights \code{a},
\code{b} and \code{c}. Comparing the codes given for those two digits is a quick way to confirm
that you have the layout the right way round before you start deriving \code{A} to \code{F}.

On the DE0-CV the displays are wired \textbf{active low}: driving a segment to \code{'0'} lights
it, and \code{'1'} turns it off. That is why the code for \code{8}, which lights every segment, is
\code{"0000000"}, and why a blank display is \code{"1111111"}.

Throughout, \code{hex(6 downto 0)} maps onto the segments \code{g}, \code{f}, \code{e}, \code{d},
\code{c}, \code{b}, \code{a}: bit 6 is \code{g} and bit 0 is \code{a}.

\exercise{\code{display}}{VHDL}
\label{c2:ex:display}
Write a module named \code{display} that drives one 7-segment display with a hexadecimal digit.

The entity must have:

\begin{booktable}{@{}p{5em}p{3.4em}p{10.8em}L@{}}
  \thead{Port} & \thead{Dir.} & \thead{Type} & \thead{Description} \\
  \midrule
  \code{number} & in & \code{std_logic_vector(3 downto 0)} & The value to show, \code{0000} through
    \code{1111}. \\
  \code{hex} & out & \code{std_logic_vector(6 downto 0)} & The segment code for that value. \\
\end{booktable}

Feeding \code{number} the value \code{0111} must therefore put the code for \code{7} on \code{hex},
and feeding it \code{1110} must put the code for \code{E} on \code{hex}.

Implement it with a process:
\begin{itemize}
  \item Put \code{number} in the sensitivity list.
  \item Use a \code{case} statement on \code{number}, one branch per value.
  \item Include a \code{when others} branch, and blank the display in it.
\end{itemize}

You are given the codes for \code{0} to \code{9}. Work out \code{A} to \code{F} yourself, from the
segment layout above:

\begin{vhdlcode}
constant DISPLAY_0  : std_logic_vector(6 downto 0) := "1000000";
constant DISPLAY_1  : std_logic_vector(6 downto 0) := "1111001";
constant DISPLAY_2  : std_logic_vector(6 downto 0) := "0100100";
constant DISPLAY_3  : std_logic_vector(6 downto 0) := "0110000";
constant DISPLAY_4  : std_logic_vector(6 downto 0) := "0011001";
constant DISPLAY_5  : std_logic_vector(6 downto 0) := "0010010";
constant DISPLAY_6  : std_logic_vector(6 downto 0) := "0000010";
constant DISPLAY_7  : std_logic_vector(6 downto 0) := "1111000";
constant DISPLAY_8  : std_logic_vector(6 downto 0) := "0000000";
constant DISPLAY_9  : std_logic_vector(6 downto 0) := "0010000";
-- Todo: add DISPLAY_A through DISPLAY_F here.
constant DISPLAY_OFF: std_logic_vector(6 downto 0) := "1111111";
\end{vhdlcode}

Two of the six are lower case on a 7-segment display, because their capital forms would be
impossible to tell from a digit: \code{b} would read as \code{8} and \code{d} would read as
\code{0}. \code{A}, \code{C}, \code{E} and \code{F} are capitals.

All sixteen codes are also in \file{display_tb.vhd}, since a self-checking testbench cannot check an
output without knowing what it should be. So this is no puzzle you can be locked out of. Derive the
six anyway: it takes a couple of minutes with the segment diagram, and reading them out of the
testbench teaches you nothing the \code{case} statement does not. The work in this exercise is the
sixteen-branch \code{case} statement, not the lookup table.

\modulefig[0.68]{lectures/L02/appendix/images/display.png}

\begin{selfcheck}
\textbf{Self-check.} Name your entity \code{display}, with the ports \code{number}
(\code{std_logic_vector(3 downto 0)}) and \code{hex} (out, \code{std_logic_vector(6 downto 0)}),
declared in that order; its testbench is in \file{exercises/display/} and checks all sixteen values
against the whole code table, naming the digit and both codes when one is wrong.
\end{selfcheck}

\exercise{\code{hex_display}}{VHDL}
\label{c2:ex:hexdisplay}
Write a module named \code{hex_display} that shows a two-digit hexadecimal number on two displays,
using two instances of your \code{display} from Exercise~\ref{c2:ex:display}.

The entity must have:

\begin{booktable}{@{}p{4.6em}p{3.4em}p{10.8em}L@{}}
  \thead{Port} & \thead{Dir.} & \thead{Type} & \thead{Description} \\
  \midrule
  \code{input} & in & \code{std_logic_vector(7 downto 0)} & Eight switches: the top four are the
    left digit, the bottom four the right. \\
  \code{hex1} & out & \code{std_logic_vector(6 downto 0)} & Segment code for the left digit. \\
  \code{hex0} & out & \code{std_logic_vector(6 downto 0)} & Segment code for the right digit. \\
\end{booktable}

When it works, the number on \code{input(7 downto 4)} appears on \code{hex1}, and the number on
\code{input(3 downto 0)} on \code{hex0}.

Build it out of submodules, following \secref{c2:sec:submodules}:
\begin{itemize}
  \item Create two instances of \code{display}, labelled \code{display1} and \code{display0}.
  \item Connect \code{input(7 downto 4)} and \code{hex1} to \code{display1}.
  \item Connect \code{input(3 downto 0)} and \code{hex0} to \code{display0}.
  \item Use a positional \code{port map}, as everywhere else in this book.
  \item Write no logic of your own in \code{hex_display}. If you catch yourself writing a
    \code{case} statement here, you are duplicating Exercise~\ref{c2:ex:display} instead of reusing
    it.
\end{itemize}

Copy your \file{display.vhd} from Exercise~\ref{c2:ex:display} into the exercise directory, since
\code{hex_display} cannot be analysed without it. This is the book's first design built from more
than one file, so name \file{display.vhd} first, before the file that instantiates it:

\begin{shell}
ghdl -a --std=93 display.vhd hex_display.vhd hex_display_tb.vhd
ghdl -e --std=93 hex_display_tb
ghdl -r --std=93 hex_display_tb --assert-level=error --stop-time=10ms
\end{shell}

\Secref{c2:sub:multifile} explains why the order matters.

\modulefig[0.68]{lectures/L02/appendix/images/hex_display.png}

\begin{selfcheck}
\textbf{Self-check.} Name your entity \code{hex_display}, with the ports \code{input}
(\code{std_logic_vector(7 downto 0)}), \code{hex1} (out) and \code{hex0} (out, both
\code{std_logic_vector(6 downto 0)}), declared in that order; its testbench is in
\file{exercises/hex_display/} and sweeps all 256 input values, checking each half independently so
that a design swapping the two digits is caught instead of passing on the values where both halves
happen to agree.
\end{selfcheck}

\exercise{One digit as a gate network}{Circuit}
\label{c2:ex:segments}
\emph{If you have time.} Realise one digit as a gate network, the way \chapref{c1:ch} and
\secref{c2:sec:kmap} had you do:
\begin{itemize}
  \item Take the four inputs and one segment, say \code{a}, and derive the Boolean equation for it
    from the 16-row truth table. A Karnaugh map is worth using here: most segments minimise
    considerably.
  \item Repeat for as many of the six remaining segments as you have patience for, and build the
    result in CircuitVerse.
  \item Then copy the whole gate network to drive a second display.
\end{itemize}

That copy is the point of the exercise. Duplicating a block of gates is exactly what instantiating
\code{display} a second time did in Exercise~\ref{c2:ex:hexdisplay}, and it is worth having done
once by hand to see what that single line of VHDL stands for.

% ------------------------------------------------------------------------------------------
\exerciseset{Naming an intermediate result}

\exercise{\code{combo_logic}, written twice}{VHDL}
\label{c2:ex:combo}
Write an entity named \code{combo_logic} that implements the Boolean function

\[ x = ab + c'd \]

\Secref{c2:sec:tovhdl} claimed that naming a subexpression with an internal signal costs nothing and
buys readability. This exercise is where you confirm that rather than take it on faith: you write
the function twice, once as a single expression and once with the term $c'd$ factored out into an
internal signal named \code{y}.

A signal declared inside an architecture is a wire. It is not a variable, it stores nothing, and
naming a subexpression with one adds neither a gate, a delay nor a flip-flop to the circuit. It
only gives part of the network a name, so that when you read it later you can see what that part is
for. Both architectures below synthesise to the same gates, and the testbench handed out passes
against either of them.

The entity must have the ports \code{a}, \code{b}, \code{c}, \code{d} (in) and \code{x} (out), all
\code{std_logic}.

\xpart{a} Implement the whole function with a single concurrent signal assignment: declare no
internal signals, and translate every Boolean operation into the corresponding VHDL operator.

\xpart{b} Rewrite the architecture with an internal signal named \code{y}, representing the term
$c'd$, and assign \code{x} the final expression. This is the same shape as \code{gate_network} in
\secref{c2:sec:tovhdl}, which names its \code{AND} term \code{ab}.

Confirm that the two architectures describe the same circuit.

\note[Hint]Two descriptions represent the same combinational circuit when they produce the same
output for every possible input combination. With four inputs there are 16 possible input
combinations: write the full truth table and verify that both architectures produce the same value
on \code{x} on every row.

\modulefig[0.52]{lectures/L02/appendix/images/combo_logic.png}

\begin{selfcheck}
\textbf{Self-check.} Name your entity \code{combo_logic}, with the ports \code{a}, \code{b},
\code{c}, \code{d} (in) and \code{x} (out), declared in that order; its testbench is in
\file{exercises/combo_logic/} and checks all 16 rows. It passes against either architecture, which
is the point of the comparison above.
\end{selfcheck}
